\documentclass[11pt,a4paper]{article}
\usepackage[utf8]{inputenc}
\usepackage[T1]{fontenc}
\usepackage[french]{babel}
\usepackage{geometry}
\usepackage{booktabs}
\usepackage{longtable}
\usepackage{array}
\usepackage{xcolor}
\usepackage{colortbl}
\usepackage{graphicx}
\usepackage{fancyhdr}
\usepackage{titlesec}
\usepackage{tikz}
\usepackage{enumitem}
\usepackage{hyperref}
\usepackage{multicol}
\usepackage{listings}
\usepackage{tcolorbox}
\usepackage{amssymb}
\usepackage{amsmath}
\usepackage{float}


\geometry{margin=2cm}
\pagestyle{fancy}
\fancyhf{}
\rhead{Atelier Machine Learning}
\lhead{Analyse Comportementale Clientèle}
\rfoot{Page \thepage}

% Couleurs personnalisées
\definecolor{headerblue}{RGB}{41, 65, 114}
\definecolor{lightgray}{RGB}{245, 245, 245}
\definecolor{accentorange}{RGB}{230, 126, 34}
\definecolor{codegreen}{RGB}{0,128,0}
\definecolor{codegray}{RGB}{128,128,128}
\definecolor{codepurple}{RGB}{128,0,128}
\definecolor{backcolour}{RGB}{248,248,248}
\definecolor{warningred}{RGB}{231, 76, 60}
\definecolor{successgreen}{RGB}{46, 204, 113}

\lstdefinestyle{mystyle}{
    backgroundcolor=\color{backcolour},
    commentstyle=\color{codegreen},
    keywordstyle=\color{codepurple},
    numberstyle=\tiny\color{codegray},
    stringstyle=\color{codegreen},
    basicstyle=\ttfamily\footnotesize,
    breakatwhitespace=false,
    breaklines=true,
    captionpos=b,
    keepspaces=true,
    numbers=left,
    numbersep=5pt,
    showspaces=false,
    showstringspaces=false,
    showtabs=false,
    tabsize=2
}
\lstset{style=mystyle}


\tcbset{
    colback=backcolour,
    colframe=headerblue,
    fonttitle=\bfseries,
    boxrule=1pt,
    arc=3mm
}


\titleformat{\section}
{\normalfont\Large\bfseries}
{\thesection}{1em}{}[{\titlerule[2pt]}]

\titleformat{\subsection}
{\normalfont\large\bfseries}
{\thesubsection}{1em}{}

% Définition de colonnes personnalisées
\newcolumntype{C}[1]{>{\centering\arraybackslash}p{#1}}
\newcolumntype{L}[1]{>{\raggedright\arraybackslash}p{#1}}

\begin{document}


\begin{titlepage}
\vspace{2cm}

\begin{center}
{\Huge\textbf{Atelier Machine Learning}}\par
\vspace{0.5cm}
{\Large\textit{Analyse Comportementale Clientèle Retail}}\par

\vspace{4cm}

{\LARGE\textbf{COMPTE RENDU}}\par
\vspace{1cm}
{\large Atelier Pratique : E-commerce de Cadeaux}\par
\vspace{2cm}
{\large \textbf{Préparé par XXXX}}
\vspace{3cm}

\rule{10cm}{2pt}

\vspace{5cm}

{\textbf{Année Univeristaire : 2025-2026}}

\end{center}
\vfill
\end{titlepage}

\tableofcontents
\newpage

%==============================================================================
\section{Introduction}
%==============================================================================

Ce rapport présente les travaux réalisés dans le cadre de l'atelier Machine Learning portant sur l'analyse comportementale de la clientèle d'une boutique e-commerce spécialisée dans les cadeaux. L'objectif principal de ce projet est de développer des solutions d'apprentissage automatique permettant de :

\begin{itemize}
    \item \textbf{Personnaliser la stratégie marketing} grâce à la segmentation client
    \item \textbf{Réduire le taux d'attrition} (churn) par la modélisation prédictive
    \item \textbf{Augmenter le chiffre d'affaires} en ciblant les clients à risque
\end{itemize}

Le jeu de données mis à disposition contient \textbf{4 372 clients} décrits par \textbf{52 variables}, incluant des métriques RFM (Récence, Fréquence, Montant), des données démographiques, des comportements d'achat et des informations de compte.

%==============================================================================
\section{Prétraitement des Données}
%==============================================================================

Le prétraitement constitue une étape cruciale pour garantir la qualité des modèles. Nous avons implémenté un pipeline complet dans le fichier \texttt{src/preprocessing.py}.

\subsection{Gestion des Valeurs Manquantes}

\subsubsection{Variable Age}
La variable \texttt{Age} présentait \textbf{1 311 valeurs manquantes} (29,99\% du dataset). Nous avons opté pour une imputation par la méthode des \textbf{K plus proches voisins (KNN)} avec les paramètres suivants :
\begin{itemize}
    \item Nombre de voisins : $k = 5$
    \item Pondération par la distance inverse
    \item Variables utilisées : Recency, Frequency, MonetaryTotal, MonetaryAvg, CustomerTenureDays, FirstPurchaseDaysAgo
\end{itemize}

Cette approche permet de prendre en compte les corrélations entre l'âge et les comportements d'achat, offrant une imputation plus réaliste qu'une simple moyenne ou médiane.

\subsubsection{Variable AvgDaysBetweenPurchases}
Cette variable contenait \textbf{79 valeurs manquantes}, imputées par la médiane (1,15 jours).

\subsection{Traitement des Valeurs Aberrantes}

Certaines variables contenaient des \textbf{valeurs sentinelles} utilisées pour encoder des données manquantes ou des erreurs :

\begin{table}[H]
\centering
\begin{tabular}{lll}
\toprule
\textbf{Variable} & \textbf{Valeurs aberrantes} & \textbf{Nombre d'occurrences} \\
\midrule
SupportTicketsCount & -1, 999 & 130 (2,97\%) \\
SatisfactionScore & -1, 0, 99 & 349 (7,98\%) \\
\bottomrule
\end{tabular}
\caption{Valeurs sentinelles identifiées et corrigées}
\end{table}

Ces valeurs ont été remplacées par la médiane de chaque variable respective, puis bornées aux plages valides (0-15 pour SupportTicketsCount, 1-5 pour SatisfactionScore).

\subsection{Ingénierie des Variables}

\subsubsection{Date d'inscription}
La variable \texttt{RegistrationDate} a été parsée et décomposée en quatre nouvelles variables :
\begin{itemize}
    \item \texttt{RegYear} : Année d'inscription
    \item \texttt{RegMonth} : Mois d'inscription
    \item \texttt{RegDay} : Jour d'inscription
    \item \texttt{RegWeekday} : Jour de la semaine (0=Lundi, 6=Dimanche)
\end{itemize}

\subsubsection{Adresse IP}
À partir de la variable \texttt{LastLoginIP}, nous avons extrait :
\begin{itemize}
    \item \texttt{IsPrivateIP} : Indicateur booléen (326 clients avec IP privée)
    \item \texttt{LoginCountry} : Pays de connexion (via la base DB-IP Lite)
\end{itemize}

La distribution des pays de connexion révèle une clientèle majoritairement américaine (1 494 clients), suivie de la Chine (362) et du Japon (219).

\subsection{Suppression des Variables Non Fiables}

La variable \texttt{NewsletterSubscribed} a été supprimée car identifiée comme non fiable dans la documentation du projet.

%==============================================================================
\section{Détection et Correction des Fuites de Données}
%==============================================================================

Une analyse approfondie des corrélations entre les variables et la cible \texttt{Churn} a révélé plusieurs cas de \textbf{fuite de données} (data leakage) compromettant la validité des modèles.

\subsection{Variables Exclues}

\begin{table}[H]
\centering
\begin{tabular}{L{3.5cm}L{9cm}}
\toprule
\textbf{Variable} & \textbf{Raison de l'exclusion} \\
\midrule
ChurnRiskCategory & Directement dérivée de la variable cible (100\% de corrélation) \\
RFMSegment & La catégorie "Dormants" correspond à 100\% de clients churned \\
CustomerType & La catégorie "Perdu" correspond à 100\% de clients churned \\
Recency & Séparation parfaite : Churn=0 si Recency $\leq$ 90 jours, Churn=1 si Recency $>$ 90 jours \\
\bottomrule
\end{tabular}
\caption{Variables exclues pour cause de fuite de données}
\end{table}

\subsection{Analyse de la Variable Recency}

L'analyse a montré que la variable \texttt{Recency} définit elle-même le churn :
\begin{itemize}
    \item Clients non churned : Recency $\in [1, 90]$ jours
    \item Clients churned : Recency $\in [91, 374]$ jours
\end{itemize}

Cette découverte est critique car elle indique que le churn est défini comme l'inactivité supérieure à 90 jours. Utiliser cette variable reviendrait à prédire le passé avec le passé.

%==============================================================================
\section{Modélisation Prédictive du Churn}
%==============================================================================

\subsection{Préparation des Données pour la Modélisation}

\subsubsection{Encodage des Variables}
\begin{itemize}
    \item \textbf{Variables numériques (35)} : Standardisation avec \texttt{StandardScaler}
    \item \textbf{Variables catégorielles (14)} : Encodage one-hot avec \texttt{OneHotEncoder}
\end{itemize}

\subsubsection{Séparation Train/Test}
Le dataset a été divisé en :
\begin{itemize}
    \item \textbf{Ensemble d'entraînement} : 3 497 échantillons (80\%)
    \item \textbf{Ensemble de test} : 875 échantillons (20\%)
\end{itemize}

La stratification a été appliquée pour maintenir la distribution de la variable cible (33,26\% de churn) dans les deux ensembles.

\subsubsection{Gestion du Déséquilibre de Classes}
Le taux de churn de 33,26\% représente un déséquilibre modéré. Nous avons appliqué la technique \textbf{SMOTE} (Synthetic Minority Over-sampling Technique) sur l'ensemble d'entraînement :
\begin{itemize}
    \item Avant SMOTE : 2 334 non-churned, 1 163 churned
    \item Après SMOTE : 2 334 non-churned, 2 334 churned
\end{itemize}

\subsection{Modèles Évalués}

Sept algorithmes de classification ont été entraînés et comparés :

\begin{table}[H]
\centering
\begin{tabular}{lccccc}
\toprule
\textbf{Modèle} & \textbf{Accuracy} & \textbf{Precision} & \textbf{Recall} & \textbf{F1} & \textbf{ROC-AUC} \\
\midrule
\textbf{XGBoost (Tuned)} & \textbf{99,1\%} & \textbf{99,0\%} & \textbf{98,3\%} & \textbf{98,6\%} & \textbf{99,97\%} \\
Logistic Regression & 98,4\% & 98,6\% & 96,6\% & 97,6\% & 99,8\% \\
Gradient Boosting & 98,3\% & 96,9\% & 97,9\% & 97,4\% & 99,9\% \\
Decision Tree & 97,3\% & 94,9\% & 96,9\% & 95,9\% & 97,2\% \\
SVM & 97,1\% & 98,9\% & 92,4\% & 95,6\% & 99,7\% \\
Random Forest & 94,9\% & 96,6\% & 87,6\% & 91,9\% & 98,7\% \\
KNN & 84,1\% & 70,9\% & 88,7\% & 78,8\% & 92,5\% \\
\bottomrule
\end{tabular}
\caption{Comparaison des performances des modèles sur l'ensemble de test}
\end{table}

\subsection{Optimisation des Hyperparamètres}

Le modèle XGBoost, identifié comme le plus performant selon le score F1, a fait l'objet d'une optimisation par \textbf{GridSearchCV} avec validation croisée à 5 plis.

\begin{tcolorbox}[title=Hyperparamètres optimaux - XGBoost, colback=white, colframe=black]
\begin{itemize}
    \item \texttt{learning\_rate} : 0,1
    \item \texttt{max\_depth} : 7
    \item \texttt{n\_estimators} : 100
    \item \texttt{subsample} : 0,8
\end{itemize}
Score F1 en validation croisée : 99,09\%
\end{tcolorbox}

\subsection{Importance des Variables}

L'analyse de l'importance des variables du modèle XGBoost révèle les facteurs les plus prédictifs du churn :

\begin{enumerate}
    \item \textbf{LoyaltyLevel\_Établi} (35,3\%) : Les clients établis sont significativement moins susceptibles de churner
    \item \textbf{FavoriteSeason\_Hiver} (20,4\%) : Les habitudes d'achat saisonnières sont fortement prédictives
    \item \textbf{PreferredMonth} (20,0\%) : Le mois d'achat préféré indique le niveau d'engagement
    \item \textbf{LoyaltyLevel\_Nouveau} (5,8\%) : Les nouveaux clients présentent des patterns distincts
    \item \textbf{CustomerTenureDays} (4,5\%) : L'ancienneté du client
\end{enumerate}

%==============================================================================
\section{Segmentation Client}
%==============================================================================

En complément de la prédiction du churn, nous avons réalisé une segmentation client basée sur l'analyse RFM pour personnaliser les stratégies marketing.

\subsection{Méthodologie}

\subsubsection{Variables RFM}
\begin{itemize}
    \item \textbf{Recency} : Nombre de jours depuis le dernier achat
    \item \textbf{Frequency} : Nombre total d'achats
    \item \textbf{Monetary} : Montant total dépensé
\end{itemize}

\subsubsection{Algorithme de Clustering}
Nous avons utilisé l'algorithme \textbf{K-Means} avec les étapes suivantes :
\begin{enumerate}
    \item Standardisation des variables RFM avec \texttt{StandardScaler}
    \item Recherche du nombre optimal de clusters via le score de silhouette
    \item Application de K-Means avec $k=3$ clusters
\end{enumerate}

\subsection{Segments Identifiés}

\begin{table}[H]
\centering
\begin{tabular}{L{3cm}ccccc}
\toprule
\textbf{Segment} & \textbf{Clients} & \textbf{Taux Churn} & \textbf{Récence moy.} & \textbf{Fréq. moy.} & \textbf{Montant moy.} \\
\midrule
À Risque & 1 109 (25,4\%) & 100\% & 246 jours & 1,85 & 460 \$ \\
Champions & 26 (0,6\%) & 0\% & 6 jours & 83,35 & 75 966 \$ \\
Clients Actifs & 3 237 (74,0\%) & 11\% & 40 jours & 5,55 & 1 796 \$ \\
\bottomrule
\end{tabular}
\caption{Profil des segments clients identifiés}
\end{table}

\subsection{Recommandations Marketing par Segment}

\begin{tcolorbox}[title=Segment Champions (Priorité : HAUTE VALEUR), colback=white, colframe=black]
\textbf{Actions recommandées :}
\begin{itemize}
    \item Offres exclusives et accès anticipé aux nouveaux produits
    \item Programme de fidélité avec avantages VIP
    \item Sollicitation d'avis et programme de parrainage
    \item Communications personnalisées de remerciement
\end{itemize}
\textbf{Canaux :} Email, Courrier direct, Gestionnaire de compte dédié
\end{tcolorbox}

\begin{tcolorbox}[title=Segment À Risque (Priorité : URGENTE), colback=white, colframe=black]
\textbf{Actions recommandées :}
\begin{itemize}
    \item Campagne de reconquête avec remises spéciales
    \item Enquête pour comprendre les raisons de l'inactivité
    \item Messages personnalisés "Vous nous manquez"
    \item Offres exclusives à durée limitée
\end{itemize}
\textbf{Canaux :} Email, SMS, Courrier direct
\end{tcolorbox}

\begin{tcolorbox}[title=Segment Clients Actifs (Priorité : MOYENNE), colback=white, colframe=black]
\textbf{Actions recommandées :}
\begin{itemize}
    \item Up-selling vers des produits de valeur supérieure
    \item Cross-selling de produits complémentaires
    \item Seuils de livraison gratuite pour augmenter le panier moyen
    \item Accès anticipé aux soldes
\end{itemize}
\textbf{Canaux :} Email, Publicités de retargeting
\end{tcolorbox}

%==============================================================================
\section{Visualisations Générées}
%==============================================================================

Le projet génère six visualisations clés stockées dans le répertoire \texttt{reports/}.

\subsection{Comparaison des Modèles (model\_comparison.png)}

Cette visualisation présente deux graphiques en barres côte à côte :
\begin{itemize}
    \item \textbf{Panneau gauche} : Métriques de classification (Accuracy, Precision, Recall, F1) pour chaque modèle
    \item \textbf{Panneau droit} : Score ROC-AUC par modèle
\end{itemize}

Elle permet d'identifier rapidement XGBoost comme le modèle le plus performant sur tous les critères.

\subsection{Courbes ROC (roc\_curves.png)}

Affiche les courbes ROC (Receiver Operating Characteristic) de tous les modèles sur un même graphique :
\begin{itemize}
    \item Axe X : Taux de faux positifs (FPR)
    \item Axe Y : Taux de vrais positifs (TPR)
    \item La ligne diagonale représente un classifieur aléatoire (AUC = 0,5)
\end{itemize}

Les courbes proches du coin supérieur gauche indiquent une meilleure performance.

\subsection{Matrice de Confusion (confusion\_matrix\_best.png)}

Présente les résultats de prédiction du modèle XGBoost sous forme de heatmap :
\begin{itemize}
    \item \textbf{Vrais Négatifs (TN)} : Clients correctement identifiés comme non-churned
    \item \textbf{Vrais Positifs (TP)} : Clients correctement identifiés comme churned
    \item \textbf{Faux Positifs (FP)} : Clients non-churned incorrectement classés (erreur de Type I)
    \item \textbf{Faux Négatifs (FN)} : Clients churned manqués (erreur de Type II)
\end{itemize}

\subsection{Importance des Variables (feature\_importance.png)}

Graphique en barres horizontales montrant les 20 variables les plus influentes du modèle XGBoost, classées par leur contribution à la réduction des erreurs de prédiction.

\subsection{Distribution des Segments (segment\_distribution.png)}

Visualisation en deux panneaux :
\begin{itemize}
    \item \textbf{Diagramme circulaire} : Répartition des clients par segment
    \item \textbf{Diagramme en barres} : Taux de churn par segment
\end{itemize}

\subsection{Heatmap des Profils RFM (segment\_heatmap.png)}

Matrice colorée montrant les profils RFM normalisés de chaque segment :
\begin{itemize}
    \item Lignes : Métriques RFM (Récence, Fréquence, Montant)
    \item Colonnes : Segments clients
    \item Valeurs : Scores normalisés entre 0 et 1
\end{itemize}

%==============================================================================
\section{Architecture du Projet}
%==============================================================================

\subsection{Structure des Répertoires}

\begin{lstlisting}[language=bash]
ml-retail/
├── data/
│   ├── raw/                    # Données brutes
│   ├── processed/              # Données nettoyées
│   └── train_test/             # Ensembles train/test
├── models/                     # Modèles entraînés
├── reports/                    # Visualisations
└── src/                        # Code source
    ├── preprocessing.py        # Pipeline de prétraitement
    ├── train_model.py          # Entraînement des modèles
    ├── segmentation.py         # Segmentation client
    ├── predict.py              # Script de prédiction
    └── utils.py                # Fonctions utilitaires
\end{lstlisting}

\subsection{Artefacts Générés}

\begin{table}[H]
\centering
\begin{tabular}{ll}
\toprule
\textbf{Fichier} & \textbf{Description} \\
\midrule
\texttt{best\_model\_xgboost\_tuned.joblib} & Modèle XGBoost optimisé \\
\texttt{preprocessor.joblib} & Pipeline de prétraitement \\
\texttt{kmeans\_segmentation.joblib} & Modèle de clustering K-Means \\
\texttt{feature\_names.json} & Noms des variables après encodage \\
\texttt{model\_results.csv} & Résultats comparatifs des modèles \\
\texttt{segment\_profiles.csv} & Profils des segments \\
\texttt{marketing\_recommendations.json} & Recommandations par segment \\
\bottomrule
\end{tabular}
\caption{Artefacts sauvegardés dans le répertoire models/}
\end{table}

%==============================================================================
\section{Conclusion}
%==============================================================================

Ce projet a permis de développer une solution complète d'analyse comportementale client comprenant :

\begin{enumerate}
    \item \textbf{Un pipeline de prétraitement robuste} gérant les valeurs manquantes (KNN pour l'âge), les valeurs aberrantes (valeurs sentinelles) et l'ingénierie de variables (extraction de features temporelles et géographiques).

    \item \textbf{Une détection rigoureuse des fuites de données} ayant conduit à l'exclusion de 4 variables (ChurnRiskCategory, RFMSegment, CustomerType, Recency) qui auraient artificiellement gonflé les performances.

    \item \textbf{Un modèle de prédiction du churn performant} (XGBoost) atteignant un score F1 de 98,6\% et un ROC-AUC de 99,97\% sur l'ensemble de test.

    \item \textbf{Une segmentation client actionnable} identifiant trois segments distincts avec des recommandations marketing personnalisées.
\end{enumerate}

\subsection{Perspectives}

Plusieurs pistes d'amélioration peuvent être envisagées :
\begin{itemize}
    \item Intégration de données temporelles pour une analyse de séries chronologiques
    \item Développement d'un tableau de bord interactif pour le suivi en temps réel
    \item Mise en production du modèle via une API REST
    \item A/B testing des recommandations marketing par segment
\end{itemize}

\end{document}
